% Options for packages loaded elsewhere
% Options for packages loaded elsewhere
\PassOptionsToPackage{unicode}{hyperref}
\PassOptionsToPackage{hyphens}{url}
\PassOptionsToPackage{dvipsnames,svgnames,x11names}{xcolor}
%
\documentclass[
  letterpaper,
  DIV=11,
  numbers=noendperiod]{scrartcl}
\usepackage{xcolor}
\usepackage[margin=1in]{geometry}
\usepackage{amsmath,amssymb}
\setcounter{secnumdepth}{-\maxdimen} % remove section numbering
\usepackage{iftex}
\ifPDFTeX
  \usepackage[T1]{fontenc}
  \usepackage[utf8]{inputenc}
  \usepackage{textcomp} % provide euro and other symbols
\else % if luatex or xetex
  \usepackage{unicode-math} % this also loads fontspec
  \defaultfontfeatures{Scale=MatchLowercase}
  \defaultfontfeatures[\rmfamily]{Ligatures=TeX,Scale=1}
\fi
\usepackage{lmodern}
\ifPDFTeX\else
  % xetex/luatex font selection
  \setmainfont[BoldFont=EB Garamond SemiBold,ItalicFont=EB Garamond
Italic,BoldItalicFont=EB Garamond SemiBold Italic]{EB Garamond Math}
  \setmathfont[]{EB Garamond Math}
\fi
% Use upquote if available, for straight quotes in verbatim environments
\IfFileExists{upquote.sty}{\usepackage{upquote}}{}
\IfFileExists{microtype.sty}{% use microtype if available
  \usepackage[]{microtype}
  \UseMicrotypeSet[protrusion]{basicmath} % disable protrusion for tt fonts
}{}
\makeatletter
\@ifundefined{KOMAClassName}{% if non-KOMA class
  \IfFileExists{parskip.sty}{%
    \usepackage{parskip}
  }{% else
    \setlength{\parindent}{0pt}
    \setlength{\parskip}{6pt plus 2pt minus 1pt}}
}{% if KOMA class
  \KOMAoptions{parskip=half}}
\makeatother
% Make \paragraph and \subparagraph free-standing
\makeatletter
\ifx\paragraph\undefined\else
  \let\oldparagraph\paragraph
  \renewcommand{\paragraph}{
    \@ifstar
      \xxxParagraphStar
      \xxxParagraphNoStar
  }
  \newcommand{\xxxParagraphStar}[1]{\oldparagraph*{#1}\mbox{}}
  \newcommand{\xxxParagraphNoStar}[1]{\oldparagraph{#1}\mbox{}}
\fi
\ifx\subparagraph\undefined\else
  \let\oldsubparagraph\subparagraph
  \renewcommand{\subparagraph}{
    \@ifstar
      \xxxSubParagraphStar
      \xxxSubParagraphNoStar
  }
  \newcommand{\xxxSubParagraphStar}[1]{\oldsubparagraph*{#1}\mbox{}}
  \newcommand{\xxxSubParagraphNoStar}[1]{\oldsubparagraph{#1}\mbox{}}
\fi
\makeatother


\usepackage{longtable,booktabs,array}
\newcounter{none} % for unnumbered tables
\usepackage{calc} % for calculating minipage widths
% Correct order of tables after \paragraph or \subparagraph
\usepackage{etoolbox}
\makeatletter
\patchcmd\longtable{\par}{\if@noskipsec\mbox{}\fi\par}{}{}
\makeatother
% Allow footnotes in longtable head/foot
\IfFileExists{footnotehyper.sty}{\usepackage{footnotehyper}}{\usepackage{footnote}}
\makesavenoteenv{longtable}
\usepackage{graphicx}
\makeatletter
\newsavebox\pandoc@box
\newcommand*\pandocbounded[1]{% scales image to fit in text height/width
  \sbox\pandoc@box{#1}%
  \Gscale@div\@tempa{\textheight}{\dimexpr\ht\pandoc@box+\dp\pandoc@box\relax}%
  \Gscale@div\@tempb{\linewidth}{\wd\pandoc@box}%
  \ifdim\@tempb\p@<\@tempa\p@\let\@tempa\@tempb\fi% select the smaller of both
  \ifdim\@tempa\p@<\p@\scalebox{\@tempa}{\usebox\pandoc@box}%
  \else\usebox{\pandoc@box}%
  \fi%
}
% Set default figure placement to htbp
\def\fps@figure{htbp}
\makeatother





\setlength{\emergencystretch}{3em} % prevent overfull lines

\providecommand{\tightlist}{%
  \setlength{\itemsep}{0pt}\setlength{\parskip}{0pt}}



 


\linespread{1.05}
\RedeclareSectionCommand[
  afterskip=0.3ex,
  beforeskip=0.8ex
  ]{subsection}
\RedeclareSectionCommand[
  afterskip=0.1ex,
  beforeskip=0.5ex
  ]{subsubsection}
\KOMAoption{captions}{tableheading}
\makeatletter
\@ifpackageloaded{tcolorbox}{}{\usepackage[skins,breakable]{tcolorbox}}
\@ifpackageloaded{fontawesome5}{}{\usepackage{fontawesome5}}
\definecolor{quarto-callout-color}{HTML}{909090}
\definecolor{quarto-callout-note-color}{HTML}{0758E5}
\definecolor{quarto-callout-important-color}{HTML}{CC1914}
\definecolor{quarto-callout-warning-color}{HTML}{EB9113}
\definecolor{quarto-callout-tip-color}{HTML}{00A047}
\definecolor{quarto-callout-caution-color}{HTML}{FC5300}
\definecolor{quarto-callout-color-frame}{HTML}{acacac}
\definecolor{quarto-callout-note-color-frame}{HTML}{4582ec}
\definecolor{quarto-callout-important-color-frame}{HTML}{d9534f}
\definecolor{quarto-callout-warning-color-frame}{HTML}{f0ad4e}
\definecolor{quarto-callout-tip-color-frame}{HTML}{02b875}
\definecolor{quarto-callout-caution-color-frame}{HTML}{fd7e14}
\makeatother
\makeatletter
\@ifpackageloaded{caption}{}{\usepackage{caption}}
\AtBeginDocument{%
\ifdefined\contentsname
  \renewcommand*\contentsname{Table of contents}
\else
  \newcommand\contentsname{Table of contents}
\fi
\ifdefined\listfigurename
  \renewcommand*\listfigurename{List of Figures}
\else
  \newcommand\listfigurename{List of Figures}
\fi
\ifdefined\listtablename
  \renewcommand*\listtablename{List of Tables}
\else
  \newcommand\listtablename{List of Tables}
\fi
\ifdefined\figurename
  \renewcommand*\figurename{Figure}
\else
  \newcommand\figurename{Figure}
\fi
\ifdefined\tablename
  \renewcommand*\tablename{Table}
\else
  \newcommand\tablename{Table}
\fi
}
\@ifpackageloaded{float}{}{\usepackage{float}}
\floatstyle{ruled}
\@ifundefined{c@chapter}{\newfloat{codelisting}{h}{lop}}{\newfloat{codelisting}{h}{lop}[chapter]}
\floatname{codelisting}{Listing}
\newcommand*\listoflistings{\listof{codelisting}{List of Listings}}
\makeatother
\makeatletter
\makeatother
\makeatletter
\@ifpackageloaded{caption}{}{\usepackage{caption}}
\@ifpackageloaded{subcaption}{}{\usepackage{subcaption}}
\makeatother
\usepackage{bookmark}
\IfFileExists{xurl.sty}{\usepackage{xurl}}{} % add URL line breaks if available
\urlstyle{same}
% fallback for those not using the hyperref driver hyperxmp:
\makeatletter
\@ifundefined{xmpquote}{\newcommand{\xmpquote}[1]{#1}}{}
\makeatother
\hypersetup{
  pdftitle={PHIL 101: Introduction to Philosophy},
  pdfauthor={Brian Weatherson},
  colorlinks=true,
  linkcolor={blue},
  filecolor={Maroon},
  citecolor={Blue},
  urlcolor={Blue},
  pdfcreator={LaTeX via pandoc}}


\title{PHIL 101: Introduction to Philosophy}
\author{Brian Weatherson}
\date{Fall 2026}
\begin{document}
\maketitle


\begin{tcolorbox}[enhanced jigsaw, arc=.35mm, bottomrule=.15mm, bottomtitle=1mm, breakable, colback=white, colbacktitle=quarto-callout-important-color!10!white, colframe=quarto-callout-important-color-frame, coltitle=black, left=2mm, leftrule=.75mm, opacityback=0, opacitybacktitle=0.6, rightrule=.15mm, title=\textcolor{quarto-callout-important-color}{\faExclamation}\hspace{0.5em}{Draft for discussion with GSIs}, titlerule=0mm, toprule=.15mm, toptitle=1mm]

This version is for the pre-term meeting, not for students. The last
section records what is still open. Delete that section, and this box,
before the student version goes out.

\end{tcolorbox}

\textbf{Instructor}: Brian Weatherson\\
\textbf{Email}: weath@umich.edu\\
\textbf{Office Hours}: TBD\\
\textbf{Lectures}: Tuesday and Thursday, 9:00--9:50, room TBD\\
\textbf{Discussion sections}: Twice weekly, 50 minutes, times and rooms
TBD\\
\textbf{Web}: canvas.umich.edu and
\url{https://bweatherson.github.io/F26-Phil101-site/}

~

\section{Key Points}\label{key-points}

\begin{itemize}
\tightlist
\item
  This course is an introduction to philosophy through a set of
  connected questions: how we ought to reason, when a belief is
  rational, what minds and persons are, and what we owe each other.
\item
  No background is assumed. Most readings are short, and all of them are
  free.
\item
  The course meets twice a week for lecture and twice a week for
  discussion section. Attendance at both matters, and a good deal of the
  assessment happens in the room.
\item
  Assessment is spread across many small pieces rather than concentrated
  in a few large ones. Nothing worth more than 3\% falls due before the
  sixth week.
\item
  There is one essay, worth 25\%, and it comes at the end of a sequence
  designed to teach you how to write it.
\end{itemize}

\section{Course Description}\label{course-description}

Philosophy starts from questions that are easy to ask and hard to
answer. This course takes a handful of them and works through what has
been said in reply, from classical India and ancient Greece to work
published in the last few years.

We begin with reasoning. What makes an argument a good one, what
separates the inferences that are safe from the ones that pay, and why
does a run of evidence never quite settle what comes next?

The second unit asks when a belief is rational. We'll start with the
question of when perceptual beliefs are rational, and then move onto
other cases. A particular focus will be beliefs formed via testimony.
This raises both theoretical questions, what makes a belief formed via
testimony rational, and practical questions, in a world like this, who
should we trust?

The middle of the course is about minds. The big question here is one of
the oldest in philosophy: what is the relationship between minds and
bodies. We'll get to that by looking again at perception, and in
particular thinking about the conscious aspects of perception, what it
\emph{feels like} to perceive the world the way we do.

The rest of the course is about value. What makes an action right, what
makes a life go well, what makes you the same person you were at five,
and what a state may and may not do to you. We end on free speech, which
sends us back to where we started, since the case for free speech is a
claim about how a group of people reasons its way to the truth.

\section{Required Materials}\label{required-materials}

There is no textbook to buy. Every required reading is free and linked
from the course site and from Canvas.

Readings are short by design. Most weeks ask for under 5,000 words.
Three days in the term are heavier, and they are marked in the schedule
below.

You will need an iclicker or the iclicker app for this course, since
reading quizzes are answered in lecture. Details will be on Canvas
before the first class.

\section{Course Requirements}\label{course-requirements}

{\def\LTcaptype{none} % do not increment counter
\begin{longtable}[]{@{}lr@{}}
\toprule\noalign{}
Component & Weight \\
\midrule\noalign{}
\endhead
\bottomrule\noalign{}
\endlastfoot
Reading quizzes in lecture, best 15 of about 20 & 15\% \\
Discussion section, including Short Answer 1 & 20\% \\
Module quizzes, four at 2.5\% each & 10\% \\
Short Answer 2 & 10\% \\
Essay & 25\% \\
Final examination & 20\% \\
\textbf{Total} & \textbf{100\%} \\
\end{longtable}
}

\textbf{Reading quizzes.} Two or three questions in lecture on that
day's reading. These should be trivial \emph{if you've done the
reading}. Quizzes run on \textbf{twenty} of the days that carry a
reading and the best \textbf{fifteen} count, so five missed days cost
you nothing and need no explanation.

\textbf{Discussion section.} Twenty per cent of the grade is earned in
section. Part of it is Short Answer 1, described below. \textbf{WE NEED
TO DISCUSS WHAT GOES HERE BEFORE DISTIBUTING TO THE STTUDENTS}

\textbf{Module quizzes.} Four short open-book multiple-choice quizzes,
taken on Canvas, one at the end of each block of the course. These are
there to help you keep up, and they are weighted accordingly.

\textbf{Short Answer 1.} Three short questions on the first block of the
course, written by hand in section in the week of Monday, October 05.
About thirty minutes. This is the first piece of writing anyone reads.

\textbf{Short Answer 2.} Five questions on Lectures 10 to 12, done at
home, due Friday 23 October. They walk you through the moves a
philosophy essay makes, and are practice for the essay.

\textbf{Essay.} About 1,200 words on a question in ethics or welfare,
due Monday 23 November. The essay shoudl have a similar structure to
Short Answer 2. We will get comments back to you on Short Answer 2 well
before this is due.

\textbf{Final examination.} This will primarily be on the last part of
the course, since there is no work to turn in from that part. It will be
a mix of short-answer and multiple-choice.

\section{Submit Your Own Work}\label{submit-your-own-work}

Anything you put your name to, and hand in, has to be \emph{your own
work}. Copying passages from a textbook, borrowing someone else's work,
paying someone to do it for you, or handing the assignment to a chatbot,
are ways of creating documents that are not your own work, and are
violations of academic integrity. Getting someone (or something) else to
translate what you wrote is also \emph{not acceptable} for this course.

In practice, we have the following two rules on \emph{everything}
written that you turn in.

\begin{enumerate}
\def\labelenumi{\arabic{enumi}.}
\tightlist
\item
  There must be an \emph{acknowledgments} section where you record all
  the help you got, from friends, family, books, computers, etc. You
  should overshare here; it's fine to get lots of help, and as long as
  the result is your own work, we aren't going to worry what help you
  got. One exception to this: if if you do use any online tools (that
  includes search engines, now that they have language models built it)
  you must keep records of everything you ask them. Do not use anyone
  else's account, or any account you cannot access the history of, for
  any course-related purpose.
\item
  For anything you write, you have to be able to answer two questions:
  \emph{What did you mean by that?} and \emph{Why did you write that?}.
  That's what we mean by something being your own work; you can explain
  and defend it.
\end{enumerate}

\textbf{NOTE FOR GSIs - THIS IS IMPORTANT AND WILL BE SOMETHING YOU HAVE
TO DO; WE SHOULD TALK ABOUT THE LOGISTICS OF THIS AND WHETHER IT CAN BE
IMPROVED}

To that end, we will frequently \textbf{audit} the turned in work. In
practice, that means we'll ask the auditee (that might be you!) to
explain, in person, what they meant by various claims, and why they
wrote it. We will audit work that we think is suspicious, and we will
also do some audits \emph{at random}. If we audit you, this does not
mean we think you cheated; it could be random. It's annoying for us, and
I'm sure annoying for you, that we have to do this, but the evidence is
that people turning in work which is not their own is a bit out of
control.

\section{Canvas}\label{canvas}

There is a Canvas site for this course at
\url{https://canvas.umich.edu}. Readings, quizzes, and announcements
live there. Check it regularly.

There is also a course website at
\url{https://bweatherson.github.io/F26-Phil101-site/} with the slides
and other material.

\pagebreak

\section{Class Schedule}\label{class-schedule}

Readings are listed under the class they are for, and should be done
before that class. Where more than one required item is listed, do the
shortest one first.

The \emph{Stanford Encyclopedia of Philosophy} (plato.stanford.edu) and
the \emph{Internet Encyclopedia of Philosophy} (iep.utm.edu) have
entries on nearly every topic in this course. Both are free and written
by specialists. Checking them is a good move whenever a reading or a
lecture leaves you stuck, whether or not a particular entry is listed
below. These entries are often long, and the listings below will
sometimes point you to the part of an entry that's most relevant.

\subsection{Week 1: Getting Started}\label{week-1-getting-started}

\subsubsection{Tuesday, September 01 (Lecture 1) ---
Introduction}\label{tuesday-september-01-lecture-1-introduction}

What the course is about, and how it works. No reading.

\subsubsection{Thursday, September 03 (Lecture 2) --- Deduction: the
logic
game}\label{thursday-september-03-lecture-2-deduction-the-logic-game}

We play a deduction game in class. No reading.

\subsection{Week 2: Risky Inference}\label{week-2-risky-inference}

\subsubsection{Tuesday, September 08 (Lecture 3) --- Induction, analogy,
and
explanation}\label{tuesday-september-08-lecture-3-induction-analogy-and-explanation}

\textbf{Required}: Arthur Conan Doyle,
\href{https://www.gutenberg.org/ebooks/1661}{``The Boscombe Valley
Mystery,'' in \emph{The Adventures of Sherlock Holmes}} (Project
Gutenberg). One passage only, about two pages: from ``Sherlock Holmes
was transformed when he was hot upon such a scent as this'' to ``There
are several other indications, but these may be enough to aid us in our
search.'' This isn't a philosophy paper, but we'll start with the
philosophical presuppositions behind it.

\textbf{Recommended}: \href{https://www.gutenberg.org/ebooks/1661}{The
rest of the story}.

\subsubsection{Thursday, September 10 (Lecture 4) ---
Grue}\label{thursday-september-10-lecture-4-grue}

\textbf{Required}: Jonathan Weisberg,
\href{https://jonathanweisberg.org/vip/grue.html}{``The Grue Paradox,''
in \emph{Odds \& Ends}, appendix C}.

\textbf{Recommended}: Sven Neth,
\href{https://philpapers.org/archive/NETGNR.pdf}{``Goodman's New Riddle
of Induction Explained in Words of One Syllable.''}

\subsection{Week 3: When Is a Belief
Rational?}\label{week-3-when-is-a-belief-rational}

\subsubsection{Tuesday, September 15 (Lecture 5) --- Perceptual
evidence, and how many sources of knowledge there
are}\label{tuesday-september-15-lecture-5-perceptual-evidence-and-how-many-sources-of-knowledge-there-are}

\textbf{Required}: Todd R. Long,
\href{https://1000wordphilosophy.com/2023/03/19/epistemic-justification/}{``Epistemic
Justification: What is Rational Belief?''} (1000-Word Philosophy).

\textbf{Required}: SEP,
\href{https://plato.stanford.edu/entries/epistemology-india/\#KnowKnowSour}{``Epistemology
in Classical Indian Philosophy,'' §1.1}.

\textbf{Recommended}: IEP,
\href{https://iep.utm.edu/foundationalism-in-epistemology/}{``Foundationalism,''
§§1--2}.

\subsubsection{Thursday, September 17 (Lecture 6) --- Memory and the
forgotten
source}\label{thursday-september-17-lecture-6-memory-and-the-forgotten-source}

\textbf{Required}: SEP,
\href{https://plato.stanford.edu/entries/memory-episprob/\#PresGene}{``Epistemological
Problems of Memory,'' §2.4}.

\textbf{Recommended}: Lisa K. Fazio et al.,
\href{https://www.apa.org/pubs/journals/features/xge-0000098.pdf}{``Knowledge
Does Not Protect Against Illusory Truth.''}

\subsection{Week 4: Believing Other
People}\label{week-4-believing-other-people}

\subsubsection{Tuesday, September 22 (Lecture 7) --- Does testimony
stand on its
own?}\label{tuesday-september-22-lecture-7-does-testimony-stand-on-its-own}

\textbf{Required}: Dhirendra Mohan Datta,
\href{https://academic.oup.com/mind/article-abstract/XXXVI/143/354/1029875}{``Testimony
as a Method of Knowledge,'' \emph{Mind} 36 (1927): 354--358}.

\textbf{Recommended}: Dan Sperber et al.,
\href{https://www.dan.sperber.fr/wp-content/uploads/2010_clement-et-al_epistemic-vigilance.pdf}{``Epistemic
Vigilance,'' \emph{Mind \& Language} 25 (2010)}.

\subsubsection{Thursday, September 24 (Lecture 8) --- Whom do you
trust?}\label{thursday-september-24-lecture-8-whom-do-you-trust}

\textbf{Required}: Elizabeth Anderson,
\href{https://joelvelasco.net/teaching/2330/democracy-public-policy-and-lay-assessments-of-scientific-testimony1.pdf}{``Democracy,
Public Policy, and Lay Assessments of Scientific Testimony,''
\emph{Episteme} (2011), §2 (pp.~145--149)}.

\textbf{Recommended}: SEP,
\href{https://plato.stanford.edu/entries/scientific-knowledge-social/\#BigSciTruAut}{``The
Social Dimensions of Scientific Knowledge,'' §2}.

\subsection{Week 5: Science, and the
Mind}\label{week-5-science-and-the-mind}

\subsubsection{Tuesday, September 29 (Lecture 9) --- Why trust science,
and how trust
breaks}\label{tuesday-september-29-lecture-9-why-trust-science-and-how-trust-breaks}

\textbf{Required}: C. Thi Nguyen,
\href{https://aeon.co/essays/why-its-as-hard-to-escape-an-echo-chamber-as-it-is-to-flee-a-cult}{``Escape
the Echo Chamber'' (\emph{Aeon})}.

\textbf{Recommended}: Naomi Oreskes,
\href{https://news.harvard.edu/gazette/story/2019/10/in-why-trust-science-naomi-oreskes-explains-why-the-process-of-proof-is-worth-trusting/}{``Defending
science in a post-fact era'' (\emph{Harvard Gazette}, 2019)}.

\emph{Module Quiz 1 (Lectures 1--9) opens after class, due Sunday.}
(NOTE FOR GSIs; should I just post the questions earlier?)

\subsubsection{Thursday, October 01 (Lecture 10) --- Perception
reconsidered}\label{thursday-october-01-lecture-10-perception-reconsidered}

\textbf{Required}: Bertrand Russell,
\href{https://www.gutenberg.org/files/5827/5827-h/5827-h.htm}{\emph{The
Problems of Philosophy}, ch.~1, ``Appearance and Reality.''}

\textbf{Recommended}: SEP,
\href{https://plato.stanford.edu/entries/perception-problem/}{``The
Problem of Perception,'' §2, ``The Argument from Illusion''}; and IEP,
\href{https://iep.utm.edu/sense-da/}{``Sense-Data''}.

\subsection{Week 6: Consciousness}\label{week-6-consciousness}

\emph{Short Answer 1 is written in section this week.}

\subsubsection{Tuesday, October 06 (Lecture 11) --- What it is
like}\label{tuesday-october-06-lecture-11-what-it-is-like}

\textbf{Required}: Thomas Nagel,
\href{https://rintintin.colorado.edu/~vancecd/phil201/Nagel.pdf}{``What
Is It Like to Be a Bat?''}. Note this is slightly more reading than most
days in class.

\subsubsection{Thursday, October 08 (Lecture 12) --- The hard problem,
and
Mary}\label{thursday-october-08-lecture-12-the-hard-problem-and-mary}

\textbf{Required}: Tufan Kıymaz,
\href{https://1000wordphilosophy.com/2019/10/05/the-knowledge-argument-against-physicalism/}{``The
Knowledge Argument Against Physicalism''} (1000-Word Philosophy).

\textbf{Required}: Frank Jackson,
\href{https://courses.physics.illinois.edu/phys419/sp2021/Jackson1986_WhatMaryDidntKnow.pdf}{``What
Mary Didn't Know''}. It's probably best to read Kıymaz first as
background for Jackson's paper.

\emph{Short Answer 2 goes out after class, on Lectures 10 to 12, due
Friday 23 October.}

\subsection{Week 7: Dualism}\label{week-7-dualism}

\subsubsection{Tuesday, October 13 (Lecture 13) --- The case
for}\label{tuesday-october-13-lecture-13-the-case-for}

\textbf{Required}: Peter Adamson,
\href{https://aeon.co/ideas/what-can-avicenna-teach-us-about-the-mind-body-problem}{``What
can Avicenna teach us about the mind-body problem?'' (\emph{Aeon
Ideas})}.

\textbf{Required}: Marc Bobro,
\href{https://1000wordphilosophy.com/2018/08/05/descartes-meditations-4-6/}{``René
Descartes' \emph{Meditations} 4--6''} (1000-Word Philosophy). The most
important part is the end, on Meditation VI.

\textbf{Recommended}: Jacob Berger,
\href{https://1000wordphilosophy.com/2024/02/03/mind-body-problem/}{``The
Mind-Body Problem,'' the ``Varieties of Dualism'' section}; and
\href{https://www.earlymoderntexts.com/assets/pdfs/descartes1641.pdf}{Descartes,
Meditation VI itself (Bennett, Early Modern Texts)}.

\subsubsection{Thursday, October 15 (Lecture 14) --- The case
against}\label{thursday-october-15-lecture-14-the-case-against}

\textbf{Required}: Jacob Berger,
\href{https://1000wordphilosophy.com/2024/02/03/mind-body-problem/}{``The
Mind-Body Problem: What Are Minds?'' (1000-Word Philosophy)}.

\textbf{Recommended}: IEP,
\href{https://iep.utm.edu/identity/}{``Identity Theory'', introduction};
and Princess Elisabeth of Bohemia and Descartes,
\href{https://www.earlymoderntexts.com/assets/pdfs/descelis.pdf}{correspondence
of 1643}.

\emph{Module Quiz 2 (Lectures 10--14) opens after class, due Sunday 1
November.}

\subsection{Week 8: Consequentialism}\label{week-8-consequentialism}

\emph{Short Answer 2 due Friday, October 23.}

\subsubsection{Tuesday, October 20 --- Fall Study Break, no
class}\label{tuesday-october-20-fall-study-break-no-class}

\subsubsection{Thursday, October 22 (Lecture 15) ---
Consequentialism}\label{thursday-october-22-lecture-15-consequentialism}

\textbf{Required}: Shane Gronholz,
\href{https://1000wordphilosophy.com/2014/05/15/consequentialism/}{``Consequentialism
and Utilitarianism''} (1000-Word Philosophy).

\textbf{Required}: J. S. Mill,
\href{https://utilitarianism.net/books/utilitarianism-john-stuart-mill/2/}{\emph{Utilitarianism}}
ch.~2, first five paragraphs (extract posted to Canvas).

\textbf{Recommended}: SEP,
\href{https://plato.stanford.edu/entries/consequentialism/}{``Consequentialism''},
§1, ``Classic Utilitarianism.''

\subsection{Week 9: Deontology and
Trolleys}\label{week-9-deontology-and-trolleys}

\subsubsection{Tuesday, October 27 (Lecture 16) ---
Deontology}\label{tuesday-october-27-lecture-16-deontology}

\textbf{Required}: Andrew Chapman,
\href{https://1000wordphilosophy.com/2014/06/09/introduction-to-deontology-kantian-ethics/}{``Deontology:
Kantian Ethics''} (1000-Word Philosophy).

\textbf{Recommended}: SEP,
\href{https://plato.stanford.edu/entries/ethics-deontological/}{``Deontological
Ethics,''} §1, ``Deontology's Foil: Consequentialism''; and Immanuel
Kant,
\href{https://www.earlymoderntexts.com/assets/pdfs/kant1785.pdf}{\emph{Groundwork},
Section I}, especially the opening on the good will and the passages on
universal law.

\subsubsection{Thursday, October 29 (Lecture 17) --- Trolley problems
and
objections}\label{thursday-october-29-lecture-17-trolley-problems-and-objections}

\textbf{Required}: Gabriel Andrade,
\href{https://1000wordphilosophy.com/2023/03/11/doctrine-of-double-effect/}{``The
Doctrine of Double Effect''} (1000-Word Philosophy).

\textbf{Required}: Judith Jarvis Thomson,
\href{https://openyls.law.yale.edu/handle/20.500.13051/16338}{``The
Trolley Problem''}, sections I and II.

\subsection{Week 10: What Makes a Life Go
Well}\label{week-10-what-makes-a-life-go-well}

\subsubsection{Tuesday, November 03 (Lecture 18) --- What is
welfare?}\label{tuesday-november-03-lecture-18-what-is-welfare}

\textbf{Required}: Richard Y. Chappell and Darius Meissner,
\href{https://www.utilitarianism.net/theories-of-wellbeing}{``Theories
of Well-Being''}, the introduction.

\textbf{Required}: Roger Crisp,
\href{https://plato.stanford.edu/entries/well-being/}{``Well-Being''},
\emph{Stanford Encyclopedia of Philosophy}, §1, ``The Concept.''

\subsubsection{Thursday, November 05 (Lecture 19) ---
Hedonism}\label{thursday-november-05-lecture-19-hedonism}

\textbf{Required}: SEP,
\href{https://plato.stanford.edu/entries/well-being/}{``Well-Being''},
§4.1, ``Hedonism.''

\textbf{Required}: Robert Nozick,
\href{https://rintintin.colorado.edu/~vancecd/phil3160/Nozick1.pdf}{``The
Experience Machine''}. (This is just a two page excerpt.)

\emph{Short Answer 2 comes back Friday, with the essay question.}
(AGAIN, IT MIGHT BE BEST TO DISTRIBUTE ESSAY QUESTIONS EARLIER)

\subsection{Week 11: Desire, and the Objective
List}\label{week-11-desire-and-the-objective-list}

\subsubsection{Tuesday, November 10 (Lecture 20) --- Desire and
preference
theories}\label{tuesday-november-10-lecture-20-desire-and-preference-theories}

\textbf{Required}: SEP,
\href{https://plato.stanford.edu/entries/well-being/}{``Well-Being,''
§4.2, ``Desire Theories.''}

\textbf{Recommended}:
\href{https://utilitarianism.net/theories-of-wellbeing/\#desire-theories}{``Theories
of Well-Being'' (utilitarianism.net), the desire-theory section}.

\subsubsection{Thursday, November 12 (Lecture 21) --- Objective lists,
and who has a
welfare}\label{thursday-november-12-lecture-21-objective-lists-and-who-has-a-welfare}

\textbf{Required}: SEP,
\href{https://plato.stanford.edu/entries/well-being/}{``Well-Being''},
§4.3, ``Objective List Theories.''

\textbf{Required}:
\href{https://1000wordphilosophy.com/2014/02/03/the-moral-status-of-animals/}{``The
Moral Status of Animals''} (1000-Word Philosophy).

\textbf{Recommended}:
\href{https://aeon.co/essays/if-ais-can-feel-pain-what-is-our-responsibility-towards-them}{``If
AIs Can Feel Pain, What Is Our Responsibility Towards Them?''
(\emph{Aeon})}.

\textbf{Recommended}: Walter Sinnott-Armstrong and Vincent Conitzer,
\href{https://www.cs.cmu.edu/~conitzer/AImoralstatuschapter.pdf}{``How
Much Moral Status Could Artificial Intelligence Ever Achieve?''}. If you
are writing an essay on this topic, this is \textbf{Required}.

\emph{Module Quiz 3 (Lectures 15--21) opens after class.}

\subsection{Week 12: The Self}\label{week-12-the-self}

\subsubsection{Tuesday, November 17 (Lecture 22) --- Personal identity:
the
cases}\label{tuesday-november-17-lecture-22-personal-identity-the-cases}

\textbf{Required}: Kristin Seemuth Whaley,
\href{https://1000wordphilosophy.com/2022/02/03/psychological-approaches-to-personal-identity/}{``Psychological
Approaches to Personal Identity''} (1000-Word Philosophy).

\textbf{Required}: Derek Parfit,
\href{https://rintintin.colorado.edu/~vancecd/phil375/Parfit.pdf}{``Divided
Minds and the Nature of Persons''}.

\textbf{Recommended}: Chad Vance,
\href{https://1000wordphilosophy.com/2014/02/10/personal-identity/}{``Personal
Identity: How We Exist Over Time''} (1000-Word Philosophy).

\subsubsection{Thursday, November 19 (Lecture 23) --- What matters in
survival}\label{thursday-november-19-lecture-23-what-matters-in-survival}

\textbf{Required}: John Locke,
\href{https://www.earlymoderntexts.com/assets/pdfs/locke1690book2_4.pdf}{\emph{Essay},
II.xxvii}, §§9, 14, and 26.

\textbf{Required}: Daniel Weltman,
\href{https://1000wordphilosophy.com/2023/02/25/no-self/}{``The Buddhist
Theory of No-Self''} (1000-Word Philosophy).

\textbf{Required}: Jessica Gordon-Roth,
\href{https://plato.stanford.edu/entries/locke-personal-identity/}{``Locke
on Personal Identity''}, \emph{Stanford Encyclopedia of Philosophy}, §1,
``Locke on Persons and Personal Identity: The Basics.''

\textbf{Recommended}:
\href{https://www.earlymoderntexts.com/assets/pdfs/locke1690book2_4.pdf}{Locke,
II.xxvii, §§9--26 entire}, for anyone who wants to see the details of
Locke's view in a more original form.

\subsection{Week 13: Law}\label{week-13-law}

\emph{Essay due Monday, November 23.}

\subsubsection{Tuesday, November 24 (Lecture 24) --- Why obey the
law?}\label{tuesday-november-24-lecture-24-why-obey-the-law}

\textbf{Required}: John Protavi,
\href{https://www.protevi.com/john/FH/PDF/Crito.pdf}{``Plato's
\emph{Crito}''}. Read this first.

\textbf{Required}: Plato,
\href{https://www.gutenberg.org/files/1657/1657-h/1657-h.htm}{\emph{Crito}}.
(Jowett translation, Project Gutenberg)

\subsubsection{Thursday, November 26 --- Thanksgiving recess, no
class}\label{thursday-november-26-thanksgiving-recess-no-class}

\subsection{Week 14: Punishment and
Speech}\label{week-14-punishment-and-speech}

\subsubsection{Tuesday, December 01 (Lecture 25) --- Theories of
punishment}\label{tuesday-december-01-lecture-25-theories-of-punishment}

\textbf{Required}: Travis Joseph Rodgers,
\href{https://1000wordphilosophy.com/2019/02/05/theories-of-punishment/}{``Theories
of Punishment''} (1000-Word Philosophy).

\textbf{Recommended}: SEP,
\href{https://plato.stanford.edu/entries/legal-punishment/}{``Legal
Punishment''}, the consequentialist and retributivist sections.

\subsubsection{Thursday, December 03 (Lecture 26) --- Free speech and
the marketplace of
ideas}\label{thursday-december-03-lecture-26-free-speech-and-the-marketplace-of-ideas}

\textbf{Required}: Mark Satta,
\href{https://1000wordphilosophy.com/2021/02/04/free-speech/}{``Free
Speech''} (1000-Word Philosophy)

\textbf{Required}: J. S. Mill and H. T. Mill,
\href{https://www.gutenberg.org/files/34901/34901-h/34901-h.htm}{\emph{On
Liberty}}, extract on Canvas.

\subsection{Week 15: Ending}\label{week-15-ending}

\subsubsection{Tuesday, December 08 (Lecture 27) --- Does the
marketplace track
truth?}\label{tuesday-december-08-lecture-27-does-the-marketplace-track-truth}

\textbf{Required}: David V. Johnson,
\href{https://aeon.co/ideas/how-do-we-pry-apart-the-true-and-compelling-from-the-false-and-toxic}{``How
do we pry apart the true and compelling from the false and toxic?''
(\emph{Aeon Ideas})}.

\emph{Module Quiz 4 (Lectures 22--27) opens after class, due Friday.}

\subsubsection{Thursday, December 10 (Lecture 28) ---
Review}\label{thursday-december-10-lecture-28-review}

No reading. Bring questions.

The final examination is in the University examination period, 14--21
December. The date and room will be posted once the University publishes
the schedule.

\pagebreak

\section{Open for the GSI meeting}\label{open-for-the-gsi-meeting}

\emph{This section is not for students. Delete it before the syllabus
goes out.}

\subsection{The section component}\label{the-section-component}

We should decide what to do with the 20\%. That's a lot for you, but on
the other hand, it gives you plenty of things to do in class, and two
classes a week is a lot to prep. So I think we should go with something
like

\begin{enumerate}
\def\labelenumi{\arabic{enumi}.}
\tightlist
\item
  Short Answer 1 \textasciitilde8\%
\item
  Presentations on readings \textasciitilde6\%
\item
  Participation \textasciitilde6\%
\end{enumerate}

But we can alter point 2, or add something else. I think it would
\emph{probably} be best to be uniform across the sections, but I'm
willing to be talked out of it.

\subsection{Consequences of the November essay
deadline}\label{consequences-of-the-november-essay-deadline}

The essay is due Monday 23 November so that marked essays go back before
the examination.

Short Answer 2 is marked across two weeks, 23 October to 6 November.
That's an important deadline to hit, because the students should use the
feedback from in in essay 1. Also, this should get \emph{more} feedback
than either short answer 1 or frankly the essay. (And definitely more
than the exam.)

Then essay marking doesn't have as hard a deadline; ideally 2 weeks
after it's turned in, but a bit of slippage here isn't the end of the
world.

The stuff at the end isn't as crucial to the grading, because it only
shows up on one quiz and an exam. That may result in some loss of
enthusiasm. Now maybe that's fine; it's the end of term. And free speech
normally gets enough enthusiasm anyway. And no matter what the grade,
students over-index on the importance of the exam. But if you're
worried, we could move some of the 20\% to in class activities after the
essay is due in.

Short Answer 1 is meant to be in class, and this could cause
accommodations complications. We should flag this early, and get on top
of any issues. The same goes for any other in class writing assignments
you want to do. Also, you might want to schedule a makeup slot -
preferably one slot across all six sections, because out of 150 we will
always have some crises.




\end{document}
